\chapter*{Abstract}

Spectral graph theory is a subfield of graph theory that studies graphs using the spectral properties of associated matrices, most notably the adjacency matrix.
For an undirected graph, the adjacency matrix is symmetric and therefore has real eigenvalues. 
The set of all eigenvalues, counted with multiplicities, is called the spectrum of the graph.
Interestingly, there are strong connections between the combinatorial properties of a graph and its spectrum, which is why spectral graph theory has found numerous applications over the past decades.

In this thesis, we give an overview of the fundamentals of spectral graph theory. 
We then present some classical results that showcase the connections to the theory of expander graphs and pseudorandomness.
In the second part of the thesis, we discuss two recent applications of spectral graph theory.
The first is a long-standing geometric problem, namely determining the maximum number of equiangular lines with a fixed angle in high dimensions. 
Jiang, Tidor, Yao, Zhang, and Zhao recently solved this problem using tools of spectral graph theory and also gained new insight into the distribution of the eigenvalues of certain graphs: They showed that the second largest eigenvalue of a connected bounded-degree graph has sublinear multiplicity.
As a second application, we analyse Huang's breakthrough proof of the sensitivity conjecture, previously a major open problem in computational complexity theory.

\bigskip
\selectlanguage{german}

\begin{abstractframe}
    \setlength{\parindent}{1em}
    \footnotesize
    \noindent \textbf{Deutsche Zusammenfassung.}
    Die spektrale Graphentheorie untersucht Graphen anhand der Eigenwerte und Eigenvektoren von assoziierten Matrizen, insbesondere der Adjazenzmatrix.
    Für einen ungerichteten Graphen ist die Adjazenzmatrix symmetrisch, und hat daher reelle Eigenwerte. 
    Die Menge aller Eigenwerte, mit Vielfachheiten gezählt, wird als Spektrum des Graphen bezeichnet.
    Interessanterweise bestehen tiefgreifende Zusammenhänge zwischen den kombinatorischen Eigenschaften eines Graphen und seinem Spektrum, weshalb die spektrale Graphentheorie in den letzten Jahrzehnten zahlreiche Anwendungen gefunden hat.
    
    Diese Arbeit gibt zunächst einen Überblick über die Grundlagen der spektralen Graphentheorie. 
    Anschließend präsentieren wir einige klassische Resultate, die u.a. den Zusammenhang mit der Theorie der Expandergraphen illustrieren.
    Im zweiten Teil der Arbeit diskutieren wir zwei aktuelle Anwendungen der spektralen Graphentheorie.
    Als erstes geht es dabei um ein altes geometrisches Problem, nämlich die Bestimmung der maximalen Anzahl von Geraden im hochdimensionalen Raum, die paarweise einen gleichen Winkel einschließen.
    Jiang, Tidor, Yao, Zhang und Zhao haben dieses Problem kürzlich mit Methoden der spektralen Graphentheorie gelöst, und zugleich neue Erkenntnisse über die Verteilung der Eigenwerte bestimmter Graphen gewonnen:
    Die Vielfachheit des zweitgrößten Eigenwerts eines zusammenhängenden Graphen mit beschränktem Maximalgrad ist sublinear.
    Als zweite Anwendung analysieren wir Huang's Beweis der Sensitivity-Vermutung, einem zuvor jahrzehntelang ungelösten Problem der Komplexitätstheorie.
\end{abstractframe}

\selectlanguage{british}
